\textbf{Время:} 01:45

\textbf{Тезис:} TPC восстанавливает пространственную траекторию по заряду, который частица оставила в газе.

\textbf{На экране:} сначала золотая линия частицы в газе; затем точки ионизации и пунктирные стрелки дрейфа; после этого считывающие площадки внизу и подписи «Площадки → (x, y)», «Время дрейфа → z» справа.

\textbf{Текст:} Мы связали ширину корреляционного пика с размером источника. Теперь посмотрим, как получить необходимые для этого импульсы из сигналов детектора. Начнём с траектории частицы.

Перед нами упрощённая схема время-проекционной камеры, или TPC. Золотая линия — путь заряженной частицы через газ. Частица ионизует атомы, и вдоль её пути остаются свободные электроны: они отмечены точками.

Пунктирные стрелки показывают дальнейшее движение этих электронов. Электрическое поле переносит их к считывающей плоскости. По расположению площадок, на которые пришёл заряд, определяем две координаты. Третью получаем по времени дрейфа: скорость электронов в газе известна, поэтому время можно перевести в расстояние до плоскости считывания.

Здесь мы измеряем время движения электронов, а не время пролёта пиона. Пион уже прошёл через газ, пока созданный им заряд ещё движется к площадкам. Собирая множество таких пространственных точек, восстанавливаем траекторию частицы. В экспериментальной физике её называют треком — это определение приведено справа.

\textbf{Переход:} Траектория восстановлена. На следующем слайде используем её кривизну, чтобы определить импульс.

\textbf{Справка:} TPC — Time Projection Chamber, время-проекционная камера. В двух половинах STAR TPC электроны дрейфуют к противоположным торцам; на слайде показан один фрагмент камеры. Восстановление координаты требует известного начала отсчёта и калиброванной скорости дрейфа. Время дрейфа электронов не является временем пролёта частицы. Авторская схема по Anderson et al., NIM A 499, 659–678 (2003), §5.1–5.2; arXiv:nucl-ex/0301015.
